\scanpage{8}
\uline{Prop} Soit $L_{\bullet}$ un $k$-module $\frac12$ simplicial. Les conditions suivantes sont équivalentes

a) les $L_n$ et $\pi_n(L_{\bullet})$ sont des $k$-modules projectifs.

\deleted{b}a'{)} Les composantes \deleted{de} $(L_{\bullet}^!)_n$ de $L_{\bullet}^!$ (le complexe de chaînes de Dold-Puppe associé à $L_{\bullet}$) et les $H_n(L_{\bullet}^!)$ sont des $k$-modules projectifs

b) $L_{\unsure{*|\bullet}}$ est isomorphe à \unsure{une|un} \deleted{somme} \inserted{produit} direct, \deleted{de modules} de la forme
\[
\left(\prod_{i\geq1}\Lambda^i\Phi_{*k}\otimes_k M_i\deleted{\unsure{d}}\right)
 \times\prod_{j\geq\deleted{\unsure{1}}0}(\Lambda^j\Psi_{*k})\otimes_k N_j
\]

\deleted{\unsure{d|a}} b') $L^!$ est isomorphe à un produit direct (de complexes de cochaînes) de la forme
\[
\prod_{i\geq1}(0\longrightarrow\underbrace{M_i
 \unsure{\mathrel{-}|\longrightarrow}M_i}_{\text{en degré }i\ \text{et }i-1}
 \longrightarrow0\longrightarrow\ldots)
 \times\prod_{j\geq0}(0\longrightarrow
 \underset{\text{en degré }j}{N_j}\longrightarrow0\longrightarrow\ldots)
\]

De plus, les $M_i$ et $N_j$ de b) (resp b') sont déterminés à isom près par \deleted{\unsure{$K$}} $L_{\bullet}$, et \deleted{ce} sont les $\overline{\mathrm{m}}$\unsure{s} (à isom près) dans b) et b').

\uline{Dém} (\deleted{\unsure{b}}a'{)}$\Longleftrightarrow$b')) est trivial et bien connu, ainsi que l'unicité à isom \inserted{can} près des $M_i,N_j$ dans b') (\uline{NB} on a $N_j\simeq H_j(L_{\bullet})$, et $M_i\simeq L_i/Z_i\{\simeq B_{i-1})\deleted{\unsure{\simeq}}$. On a a)$\Longleftrightarrow$a') \inserted{et b)$\Longleftrightarrow$b')} par le \unsure{tapis} de Dold-Puppe. D'où l'équivalence des 4 conditions, et la dernière assertion est contenue.
% archive sheet number: 7
